\documentclass[11pt]{article}
\usepackage[a4paper,margin=22mm]{geometry}
\usepackage[no-math]{fontspec}
\setmainfont{DejaVu Serif}
\usepackage{amsmath,amssymb,amsthm,mathtools,graphicx,xcolor,xurl,xspace,enumerate}
\usepackage[shortlabels]{enumitem}
\setlength{\emergencystretch}{4em}
\providecommand{\smaller}{\small}
\providecommand{\loccit}{\textup{loc.\ cit.}}
\providecommand{\locscit}{\textup{locc.\ cit.}}
\providecommand{\ibid}{\textup{ibid.}}
\providecommand{\cf}{\textup{cf.}}
\providecommand{\resp}{\textup{resp.}}
\providecommand{\eg}{\textup{e.g.}}
\providecommand{\ie}{\textup{i.e.}}
\providecommand{\vs}{\textup{vs.}}
%
% retained author preamble lines below (source lines 1-150, publisher class and
% publisher meta commands removed); tex of the card body is delivered verbatim.
%

\usepackage[scr=rsfs,cal=euler]{mathalfa}
\usepackage[normalem]{ulem}
\hyphenation{alter-na-ti-vely appear-ing applied deduce direct exact explic-it-ly injec-tion insti-tute inter-chang-ing irre-ducible redu-cible under-lying}

\def\@settitle{%
\vspace*{-8mm}
\raggedleft\includegraphics[scale=.5]{titre-jep}
\vtop to 50 mm{%
\parindent=0pt
{\abstractfont\article@logo\par}
\medskip
\hrule
\vfil
\begin{center}
\def\baselinestretch{1.2}\large\vfil
{\didottitraille\MakeUppercase\@title\par}
\vfil\vfil
\begin{minipage}{.8\textwidth}\centering
\ifx\@empty\smfbyname\else
{\smf@byfont\smfbyname\ifsmf@byauthor\enspace\else\ \fi}%
\fi {\smf@authorfont \edef\smfandname{{\noexpand\smf@andfont
\smfandname}} \andify\authors\authors\par}
\end{minipage}
\vfil \vrule height .4pt width .3\textwidth \vfil
\end{center}}%
\par\enlargethispage{.5\baselineskip}%
}
\def\@setthanks{\def\thanks##1{\par##1\@addpunct{{\upshape.}}}\vspace*{-5pt}\thankses}




\newcommand\mto{\mathchoice{\longmapsto}{\mapsto}{\mapsto}{\mapsto}}
\newcommand\hto{\mathrel{\lhook\joinrel\to}}
\newcommand{\onto}{\mathchoice{\to\hspace*{-5mm}\to}{\mathrel{\hbox{$\to\hspace*{-2.85mm}\to$}}}{\to\hspace*{-3mm}\to}{\to\hspace*{-3mm}\to}}
\newcommand\from{\mathchoice{\longleftarrow}{\leftarrow}{\leftarrow}{\leftarrow}}
\def\hfrom{\mathrel{\from\joinrel\rhook}}
\newcommand{\onfrom}{\mathchoice{\from\hspace*{-5mm}\from}{\from\hspace*{-2.85mm}\from}{\from\hspace*{-3mm}\from}{\from\hspace*{-3mm}\from}}
\newcommand{\To}[1]{\mathchoice{\xrightarrow{\textstyle\kern4pt#1\kern3pt}}{\stackrel{#1}{\longrightarrow}}{}{}}
\newcommand{\quand}{\quad\text{and}\quad}
\newcommand{\sbullet}{{\scriptscriptstyle\bullet}}

\newcommand{\Psfrac}[2]{(\sfrac{#1}{#2})}
\newcommand{\psfrac}[2]{\sfrac{(#1)}{#2}}
\newcommand{\spfrac}[2]{\sfrac{#1}{(#2)}}
\newcommand{\pspfrac}[2]{\sfrac{(#1)}{(#2)}}

\let\backslash\setminus
\let\wt\widetilde
\let\ovl\overline
\let\un\underline

\usepackage[all,cmtip]{xy}\let\labelstyle\textstyle
\newdir^{ (}{{}*!/-5pt/\dir^{(}}

\renewcommand{\theenumi}{\roman{enumi}}

\newcommand{\defeq}{\stackrel{\textrm{\tiny{\upshape{def}}}}{=}}

\newcommand{\F}{\mathbb{F}}
\newcommand{\Q}{\mathbb{Q}}
\newcommand{\Z}{{\mathbb Z}}

\newcommand{\cI}{\mathcal{I}}
\newcommand{\cJ}{{\mathcal{J}}}
\newcommand{\cO}{\mathcal{O}}
\newcommand{\cS}{\mathcal{S}}
\newcommand{\cW}{\mathcal{W}}
\newcommand{\cZ}{\mathcal{Z}}

\newcommand{\D}{\mathscr D}
\renewcommand{\P}{\mathscr P}

\newcommand{\EE}{\mathrm{E}}
\newcommand{\sgn}{\mathrm{sgn}}
\renewcommand{\ss}{\mathrm{ss}}

\newcommand{\fa}{\mathfrak{a}}
\newcommand{\fb}{\mathfrak{b}}
\newcommand{\fm}{\mathfrak{m}}\let\m\fm
%\newcommand{\n}{\mathfrak{n}}

\newcommand{\Fp}{\mathbb{F}_{p}}
\newcommand{\Fq}{\mathbb{F}_{q}}
\newcommand{\Qp}{\mathbb{Q}_{p}}
\newcommand{\rbar}{\overline{r}}
\newcommand{\rhobar}{\overline{\rho}}
\newcommand{\brho}{\overline{\rho}}

\newcommand{\s}{^\times}
\newcommand{\ppar}[1]{(\mkern-3mu(#1)\mkern-3mu)}
\newcommand{\bbra}[1]{[\mkern-3mu[#1]\mkern-3mu]}

\DeclareMathOperator{\Ann}{Ann}
\DeclareMathOperator{\coker}{coker}
\DeclareMathOperator{\cosoc}{cosoc}
\DeclareMathOperator{\Ext}{Ext}
\DeclareMathOperator{\Gal}{Gal}
\DeclareMathOperator{\GL}{GL}
\DeclareMathOperator{\gr}{gr}
\DeclareMathOperator{\Hom}{Hom}
\DeclareMathOperator{\HOM}{HOM}
\DeclareMathOperator{\im}{im}
\DeclareMathOperator{\Ind}{Ind}
\DeclareMathOperator{\Inj}{Inj}
\DeclareMathOperator{\JH}{\mathrm{JH}}
\DeclareMathOperator{\PGL}{PGL}
\DeclareMathOperator{\Proj}{Proj}
\DeclareMathOperator{\rad}{rad}
\DeclareMathOperator{\soc}{soc}
\DeclareMathOperator{\Tor}{Tor}

\let\ve\varepsilon
\let\vp\varphi

\newcommand{\smatr}[4]{\bigl(\begin{smallmatrix} {#1}& {#2}\\ {#3}&{#4}\end{smallmatrix}\bigr)}
\newcommand{\xto}[1][]{\To{#1}}
\newcommand{\simto}{\xrightarrow{\,\sim\,}}\let\congto\simto

\theoremstyle{plain}
\newtheorem{prop}{Proposition}[subsection]
\newtheorem{lem}[prop]{Lemma}
\newtheorem{thm}[prop]{Theorem}
\newtheorem{cor}[prop]{Corollary}
\newtheorem{prop1}{Proposition}[section]
\newtheorem{lem1}[prop1]{Lemma}
\newtheorem{thm1}[prop1]{Theorem}
\newtheorem{cor1}[prop1]{Corollary}

\theoremstyle{definition}
\newtheorem{rem}[prop]{Remark}
\newtheorem{rem1}[prop1]{Remark}

\usepackage{cleveref}
\crefname{originaltheorem}{original theorem}{original theorems}
\Crefname{originaltheorem}{Original theorem}{Original theorems}
\makeatletter
\newlabel{thm:pi2-K1}{{6.18}{source}{}{sourcepointer-1}{}}
\newlabel{thm:pi2-K1@cref}{{[originaltheorem][0][]6.18}{[1][0][]source}{}{}{}}
\makeatother
\crefname{originalcorollary}{original corollary}{original corollarys}
\Crefname{originalcorollary}{Original corollary}{Original corollarys}
\makeatletter
\newlabel{cor:subquot-K-soc}{{6.15}{source}{}{sourcepointer-2}{}}
\newlabel{cor:subquot-K-soc@cref}{{[originalcorollary][0][]6.15}{[1][0][]source}{}{}{}}
\makeatother
\title{The K1-invariants of a subquotient as a fibre product of truncations}
\author{Christophe Breuil; Florian Herzig; Yongquan Hu; Stefano Morra; Benjamin Schraen}
\date{doi:10.5802/jep.346; CC BY 4.0}
\begin{document}
\maketitle
\noindent\textit{Editorial note.} All delivered mathematical text is the author's own, copied from the pinned publisher source of this unit; no statement or proof step was rewritten, reordered or removed. The counted statement is the labelled corollary environment at source lines 3588-3596 and the counted proof is the proof environment at source lines 3598-3621.\par\medskip
\noindent\textit{Editorial note.} The preamble of the delivered fragment is generated for this unit: the publisher class is replaced by article plus the ams packages and the author's own package, macro and theorem-declaration lines are retained in source order together with their continuation lines; publisher meta commands are removed.\par\medskip
\noindent\textit{Editorial note.} The two results invoked by the counted text are declared as non-load-bearing external pointers with their printed numbers rather than restated.\par\medskip
\noindent\textit{Allowed standard prerequisites (editorial).}\begin{itemize}
\item Identification of the natural map between coinvariants. The author's theorem (printed number 6.18). Identifies the kernel computation with a map between explicitly described modules, which is the object of the proof.
\item The socle of a subquotient. The author's earlier corollary (printed number 6.15). Supplies the description of the socle used in the reduction to the case treated by the preceding results.
\item Truncations of the fixed module and their fibre products. The author's construction of the modules attached to a parameter in the same article. Provide the modules whose fibre product appears in the statement and the filtration with respect to which the proof argues.
\end{itemize}\medskip
\ifdefined\hypertarget\hypertarget{sourcepointer-1}{}\fi
\noindent\textit{Source reference.} Original theorem 6.18: sources/jep-2026-346-on-the-constituents-of-the-mod-cohomology-of-shimura/pinned-original-source.tex, source lines 2720--2726. Prerequisite pointer invoked by the counted proof: it identifies the natural map between the coinvariant spaces with a map between the displayed modules.\par\medskip
\ifdefined\hypertarget\hypertarget{sourcepointer-2}{}\fi
\noindent\textit{Source reference.} Original corollary 6.15: sources/jep-2026-346-on-the-constituents-of-the-mod-cohomology-of-shimura/pinned-original-source.tex, source lines 2586--2595. Prerequisite pointer invoked by the counted proof for the description of the socle of a subquotient.\par\medskip
\begin{cor}\label{cor:K1-subquot}
{Assume that $\brho$ is $\max\{9,2f+3\}$-generic.}
Suppose $\pi' = \pi_1'/\pi_1$ is any nonzero subquotient, where $\pi_1 \subsetneq \pi_1' \subset \pi$.
Let $i_0 \defeq i_0(\pi_1)$, $i_0' \defeq i_0(\pi_1')$, so~$-1 \le i_0 < i_0' \le f$.
Then\vspace*{-3pt}
\begin{equation*}
\pi'^{K_1} \cong D_0(\brho^\ss)_{i_0+1} \oplus_{D_0(\brho)_{i_0+1}} (D_0(\brho)_{\le i_0'}/D_0(\brho)_{\le i_0}).
\end{equation*}
\end{cor}

\begin{proof}
Note that $\pi'^{K_1}$ is the kernel of the natural map $(\pi/\pi_1)^{K_1} \to (\pi/\pi_1')^{K_1}$.
Let us write again\vspace*{-3pt}
\[
D_{i_0} \defeq D_0(\brho^\ss)_{i_0+1} \oplus_{D_0(\brho)_{i_0+1}} (D_0(\brho)/D_0(\brho)_{\le i_0}).
\]
By~Theorem~\ref{thm:pi2-K1} the above natural map is identified with a map
\begin{equation*}
\theta: D_{i_0} \to D_{i'_0}.
\end{equation*}
Let $\theta_0 \defeq \theta|_{D_0(\brho)/D_0(\brho)_{\le i_0}}$. We~have\vspace*{-3pt}
\[
\soc_{\GL_2(\cO_K)}(\pi/\pi_1') \subset \{ \sigma \in W(\brho^\ss): \ell(\sigma) \ge i_0'+1 \}
\]
by Corollary~\ref{cor:subquot-K-soc}, and $\JH(D_0(\brho^\ss)_{i_0+1}) \cap W(\brho^\ss)$ is disjoint from that set since $i_0 < i_0'$, so~$D_0(\brho^\ss)_{i_0+1} \subset \ker(\theta)$ and hence\vspace*{-3pt}
\[
\ker(\theta) = D_0(\brho^\ss)_{i_0+1} \oplus_{D_0(\brho)_{i_0+1}} \ker(\theta_0).
\]
On~the other hand, by comparison with $\pi^{K_1} = D_0(\brho)$ we see that $\theta_0$ is the natural surjection\vspace*{-3pt}
\[
D_0(\brho)/D_0(\brho)_{\le i_0} \onto D_0(\brho)/D_0(\brho)_{\le i'_0}.
\]
The result follows.
\end{proof}

\end{document}
